% Figure: the free body of an infinitesimal band element subtending angle
% d-theta on the drum. The element carries tension T on one face and
% T + dT on the other, a normal reaction dN from the drum, and a friction
% force mu dN opposing impending slip. Integrating the element equation
% around the wrap gives the capstan (Euler-Eytelwein) relation.
\begin{figure}[htbp]
\centering
\begin{tikzpicture}[>=latex, x=1.0cm, y=1.0cm]
  \def\R{3.0}
  \coordinate (Ctr) at (0,0);
  % drum surface arc under the element
  \draw[engblue, very thick] (62:\R) arc (62:118:\R);
  \fill[engblack] (Ctr) circle (1.2pt);
  \node[font=\scriptsize, anchor=north] at (Ctr) {$O$};
  % radii to the two faces of the element
  \draw[enggray, thin] (Ctr) -- (118:\R);
  \draw[enggray, thin] (Ctr) -- (62:\R);
  % the element subtends d-theta at the centre
  \draw[->, enggray, thin] (78:0.9) arc (78:102:0.9);
  \node[enggray, font=\scriptsize] at (90:1.2) {$d\theta$};
  % band element drawn slightly outside the drum
  \def\Rb{3.18}
  \draw[engamber, line width=2.2pt] (62:\Rb) arc (62:118:\Rb);
  % tension T at the lower-right face, tangent, leaving the element
  \coordinate (Pr) at (62:\Rb);
  \draw[->, engred, very thick] (Pr) -- ($ (Pr) + (0.95,-0.50) $);
  \node[engred, font=\scriptsize, anchor=west] at ($ (Pr) + (0.95,-0.50) $) {$T$};
  % tension T + dT at the upper-left face, tangent, leaving the element
  \coordinate (Pl) at (118:\Rb);
  \draw[->, engred, very thick] (Pl) -- ($ (Pl) + (-0.95,-0.50) $);
  \node[engred, font=\scriptsize, anchor=east] at ($ (Pl) + (-0.95,-0.50) $) {$T+dT$};
  % normal reaction dN, radially outward at the element midpoint
  \coordinate (Pm) at (90:\Rb);
  \draw[->, engblue, very thick] (Pm) -- ($ (Pm) + (0,0.95) $);
  \node[engblue, font=\scriptsize, anchor=south] at ($ (Pm) + (0,0.95) $) {$dN$};
  % friction mu dN tangent at the midpoint, opposing impending slip (toward T side)
  \draw[->, englime!70!black, very thick] (Pm) -- ($ (Pm) + (0.85,0.0) $);
  \node[englime!55!black, font=\scriptsize, anchor=south west] at ($ (Pm) + (0.55,0.05) $) {$\mu\,dN$};
\end{tikzpicture}
\caption[Free body of a band element on a drum]{The free body of an
infinitesimal band element wrapping the drum through a central angle
$d\theta$. The two tangent faces carry the tension $T$ on the slack-going
side and $T+dT$ on the tight-going side; the drum pushes out on the
element with a normal reaction $dN$ and drags it tangentially with the
friction $\mu\,dN$ at impending slip. Resolving the element's equilibrium
along the normal and the tangent, and dropping the second-order products,
gives $dN = T\,d\theta$ and $dT = \mu\,dN$, so $dT/T = \mu\,d\theta$.
Integrating around the wrap angle $\beta$ yields the capstan relation
$T_1/T_2 = e^{\mu\beta}$: the tension ratio grows exponentially with the
wrap, the result the band brake and the moored line both exploit.}
\label{fig:vol03ch02:capstan-element}
\end{figure}
